\subsection{有理函数的延拓}

若把第 8 号的结果应用于维数为 1 的向量空间 \(\mathbf{V} = \mathbf{K}_s\) ，则可看出存在一个典范单射 \(\phi\) ，把 \(\mathbf{K}_s\) 映入射影直线 \(\mathbf{P}_1(\mathbf{K}) = \mathbf{P}(\mathbf{K}_s \times \mathbf{K}_s)\) ；对所有 \(\xi \in \mathbf{K}\) ， \(\phi(\xi)\) 是这样的点：其齐次坐标 \((1, \xi)\) 是关于典范基（ \(\S 1\) ，第 11 号）——即 \(\mathbf{K}_s \times \mathbf{K}_s\) 的典范基——取的。 \(\phi(\mathbf{K})\) 在 \(\mathbf{P}_1(\mathbf{K})\) 中的补集由单独一个点组成，它关于上述基的齐次坐标为 \((0, 1)\) ；它称为“无穷远点”。 \(\mathbf{P}_1(\mathbf{K})\) 也称为与 \(\mathbf{K}\) 相伴的射影域，记作 \(\tilde{\mathbf{K}}\) ， \(\tilde{\mathbf{K}}\) 中的无穷远点记作 \(\infty\) 。

特别考虑 K 为交换域的情形，并设 \(f \in \mathbf{K}(\mathbf{X})\) 是 \(\mathbf{K}\) 上一个未定元的有理函数（IV，§4）；若 \(f \neq 0\) ，则存在唯一的表达式 \(f = \alpha p/q\) ，其中 \(\alpha \in \mathbf{K}^*\) ， \(p\) 与 \(q\) 是两个互素的首一多项式（VII，§1）；设 \(m\) 与 \(n\) 分别是它们的次数，并设 \(r = \sup(m, n)\) 。我们记

\[
p _ {1} (\mathrm{T}, \mathrm{X}) = \mathrm{T} ^ {r} p (\mathrm{X} / \mathrm{T}), \quad q _ {1} (\mathrm{T}, \mathrm{X}) = \mathrm{T} ^ {r} q (\mathrm{X} / \mathrm{T});
\]

\(p_{1}\) 与 \(q_{1}\) 是 K 上两个次数为 r 的齐次多项式，满足 \(p(\mathbf{X}) = p_{1}(1, \mathbf{X})\) , \(q(\mathbf{X}) = q_{1}(1, \mathbf{X})\) 。因此，对每个元素 \(\xi \in K\) （它不是 \(q(\mathbf{X})\) 的零点）， \(f(\xi) = \alpha p(\xi)/q(\xi)\) 有定义，且我们可以写出

\[
f (\xi) = \alpha p _ {1} (1, \xi) / q _ {1} (1, \xi) = \alpha p _ {1} (\lambda , \lambda \xi) / q _ {1} (\lambda , \lambda \xi)
\]

对所有 \(\lambda \neq 0\) （在 K 中）成立。于是考虑映射

\[
(\eta , \xi) \mapsto (q _ {1} (\eta , \xi), p _ {1} (\eta , \xi))
\]

（ \(K^{2}\) 到自身的映射）；它与等价关系 \(\Delta(\mathbf{K}^{2})\) 相容，因而在过渡到商时定义了一个映射 \(\tilde{f}\) （ \(\tilde{K}\) 到自身），它在有理函数有定义的点处与 \(\xi\mapsto f(\xi)\) 一致；按语言滥用的说法，称 \(\tilde{f}\) 为 f 到 \(\tilde{K}\) 的典范延拓。

例如，若 \(f = 1 / \mathbf{X}\) ，则 \(\tilde{f}(0) = \infty\) 且 \(\tilde{f}(\infty) = 0\) ；若

\[
f = (a \mathbf {X} + b) / (c \mathbf {X} + d)
\]

其中 \(ad - bc \neq 0\) ，则 \(\tilde{f}(-d/c) = \infty\) ， \(\tilde{f}(\infty) = a/c\) （若 \(c \neq 0\) ）， \(\tilde{f}(\infty) = \infty\) （若 \(c = 0\) ）。若 \(f = a_0\mathbf{X}^n + \cdots + a_n\) 是一个次数为 \(n > 0\) 的多项式，则 \(\tilde{f}(\infty) = \infty\) 。

